% =====================================================================
%  ch04.tex —— 第 4 章 K3 曲面上凝聚层的导出范畴 · §4.1 与 §4.2
%  原书 p.567–572（自同页 “CHAPTER 4” 起，止于第 5 章 “Abelian varieties”
%  标题之前；p.567 上半页的式 (46) 与 §3.5 证明收尾见 ch03d.tex）
% =====================================================================
% [原文 p.567]
\chapter{K3 曲面上凝聚层的导出范畴}
\section{K3 曲面与 Mukai 格}
\takeaway{本节记下 K3 曲面的基本事实：秩 22 的偶上同调格、Néron--Severi 格与超越闭链格 $T_S$；再借 Todd 类的平方根把陈特征改成 Mukai 向量，得到带 Euler 配对的 Mukai 格，并证明全忠实函子所对应的闭链给出 Mukai 格之间的霍奇等距。}

本章通篇讨论复数域上 K3 曲面上凝聚层的导出范畴。我们所关心、并在本章中予以回答的主要问题如下：两个不同的 K3 曲面何时有等价的凝聚层范畴？一如既往，我们把导出范畴看作三角范畴，而把等价理解为三角范畴之间的等价。

回忆一下：对于典范类或反典范类丰富的光滑射影簇，存在一个从它的凝聚层导出范畴还原出该簇本身的程序（见定理 2.1.3）。

% [原文 p.568]
然而，对其他类型的簇，这是一个非平凡的问题；对于典范类平凡的簇，这个问题尤其有趣。

我们先列出关于 K3 曲面所需要的主要事实。回忆 K3 曲面是指满足 $K_S=0$ 与 $H^1(S,\Z)=0$ 的光滑紧合代数曲面 $S$。这类曲面实际上是单连通的。可以证明第二上同调 $H^2(S,\Z)$ 无挠，并且关于相交型构成一个秩 22 的偶格。此外，由 Noether 公式可得 $p_g(S)=1$ 与 $h^{1,1}(S)=20$。

\yizhu{这组数值与上一句的格论描述互为表里：$H^2(S,\Z)$ 上相交型的符号差是 $(3,19)$，故秩为 22；Noether 公式给出 $\chi(\cO_S)=2$，与 $c_2(S)=24$ 一起决定 $p_g=1$、$h^{1,1}=20$。原书只用一句带过，读者若初次接触 K3 曲面，宜把这三个数当作同一批事实的不同写法。【来源：原书 p.568；Barth--Hulek--Peters--Van de Ven《Compact Complex Surfaces》中 K3 曲面一章（A 级）】}

K3 曲面的主要不变量之一是它的 Néron--Severi 群 $\mathrm{NS}(S)\subset H^2(S,\Z)$，在此情形它与 Picard 群 $\Pic(S)$ 重合。$\mathrm{NS}(S)$ 的秩 $\leqslant h^{1,1}=20$。记 $T_S$ 为超越闭链格；按定义，它是 Néron--Severi 格 $\mathrm{NS}(S)$ 在第二上同调 $H^2(S,\Z)$ 中的正交补。

记 $\td_S$ 为曲面 $S$ 的 Todd 类；这个类是 $H^*(S,\Q)$ 中形如 $1+2w$ 的元素，其中 $1\in H^0(S,\Z)$ 是上同调环 $H^*(S,\Z)$ 的幺元，$w\in H^4(S,\Z)$ 是 $S$ 的基本上闭链。我们取正的平方根 $\sqrt{\td_S}=1+w$；对任何 K3 曲面，它都属于整上同调环 $H^*(S,\Z)$。

对 $S$ 上任何凝聚层都可引入陈特征，并按加性把它延拓到整个凝聚层导出范畴。若 $F$ 是 $\Db(S)$ 的对象，则其陈特征
\begin{equation*}
\ch(F)=r(F)+c_1(F)+\frac{1}{2}\bigl(c_1^2-2c_2\bigr)
\end{equation*}
属于整上同调 $H^*(S,\Z)$。

对任何对象 $F$，我们定义元素
\begin{equation*}
v(F)=\ch(F)\cdot\sqrt{\td_S}\in H^*(S,\Z),
\end{equation*}
并称之为与 $F$ 相伴的向量（或 $F$ 的 Mukai 向量）。

我们按下述规则在上同调格 $H^*(S,\Z)$ 上定义一个对称双线性型：
\begin{equation*}
(u,u')=r\cdot s'+s\cdot r'-\alpha\cdot\alpha'\in H^4(S,\Z)\cong\Z,
\end{equation*}
其中 $u=(r,\alpha,s)$、$u'=(r',\alpha',s')$ 取遍 $H^0(S,\Z)\oplus H^2(S,\Z)\oplus H^4(S,\Z)$。上同调格 $H^*(S,\Z)$ 连同这个双线性型 $(\,\cdot\,,\,\cdot\,)$ 称为 \emph{Mukai 格}，记作 $\tilde H(S,\Z)$。注意在 $H^2$ 上，双线性型 $(\,\cdot\,,\,\cdot\,)$ 与通常的相交型相差一个负号。于是 Mukai 格 $\tilde H(S,\Z)$ 同构于格 $\mathrm{U}\perp -H^2(S,\Z)$，其中 $\mathrm{U}$ 是双曲格 $\begin{pmatrix}0&1\\1&0\end{pmatrix}$，而 $\perp$ 表示正交直和。

% [Tutor Note: 文字层把 Mukai 格记号的波浪线整体丢失（抽出为 “H(S, Z)”），页图作 $\tilde H(S,\Z)$；本译按页图补出。又原书该记号用直立罗马体 H，本稿为与全书的 $H^*(S,\Z)$ 一致改用斜体 $H$ 加波浪线，差异登记在报告 §5。]

\yizhu{配对中 $-\alpha\cdot\alpha'$ 的负号来自 Serre 对偶：$\Ext^2$ 的维数由 $\Hom$ 的对偶给出，于是 $\chi$ 在中间次数上带负号。$\mathrm{U}$ 因子对应 $H^0\oplus H^4$ 上的双曲配对，正是它使两类初等等距（线丛陈特征的平移与 $r$、$s$ 的互换，见 p.571）成为 Mukai 格的自等距，后文归约 $r>1$ 全靠这一条。【来源：原书 p.568 与 p.571 的 $\varphi_m$ 公式（A 级）】}

对任何两个对象 $F$ 与 $G$，配对 $(v(F),v(G))$ 按定义是元素 $\ch(F)^{\vee}\cdot\ch(G)\cdot\td_S$ 在 $H^4$ 中的分量。于是由 Grothendieck--Riemann--Roch 定理有等式
\begin{equation*}
(v(F),v(G))=\chi(F,G):=\sum_i(-1)^i\dim\Ext^i(F,G).
\end{equation*}

格 $\tilde H(S,\Z)$ 与 $T_S$ 带有自然的霍奇结构。这里所谓霍奇结构，指的是空间 $\tilde H(S,\C)$ 与 $T_S\otimes\C$ 中固定了一个一维子空间 $H^{2,0}(S)$。

% [原文 p.569]
\begin{definition}{4.1.1}
设 $S_1$ 与 $S_2$ 是两个 K3 曲面。若两格之间存在一个把一维子空间 $H^{2,0}(S_1)$ 送到 $H^{2,0}(S_2)$ 的等距，则称 $S_1$ 与 $S_2$ 的 Mukai 格（或它们的超越闭链格）\emph{霍奇等距}。
\end{definition}

设 $\cE$ 是积 $S_1\times S_2$ 的导出范畴中任意一个对象。考虑积 $S_1\times S_2$ 上的代数闭链
\begin{equation}
Z_{\cE}:=p^*\sqrt{\td_{S_1}}\cdot\ch(\cE)\cdot\pi^*\sqrt{\td_{S_2}},\tag{47}
\end{equation}
其中 $p$ 与 $\pi$ 是下列图中的投影：
\diag{S_1\times S_2 \arrow[r,"\pi"] \arrow[d,"p"'] \& S_2\\
S_1 \&}

对 K3 曲面而言，闭链 $Z_{\cE}$ 先验是有理的，实际上却是整的。

\begin{lemma}[{{[}31{]}}]{4.1.2}
对任何对象 $\cE\in\Db(S_1\times S_2)$，陈特征 $\ch(\cE)$ 与闭链 $Z_{\cE}$ 都是整的，也就是说，它们属于 $H^*(S_1\times S_2,\Z)$。
\end{lemma}

于是闭链 $Z_{\cE}$ 定义了从 $S_1$ 的整上同调格到 $S_2$ 的整上同调格的一个映射：
\begin{equation}
\begin{array}{rccc}
f_{Z_{\cE}}\colon & H^*(S_1,\Z) & \longrightarrow & H^*(S_2,\Z)\\[2pt]
& \cup & & \cup\\[2pt]
& \alpha & \longmapsto & \pi_*(Z_{\cE}\cdot p^*(\alpha))
\end{array}\tag{48}
\end{equation}

下面的命题与 [31] 的定理 4.9 类似。

\begin{proposition}{4.1.3}
设 $\cE$ 是一个对象。若函子 $\Phi_{\cE}:\Db(S_1)\to\Db(S_2)$ 全忠实，则
\begin{enumerate}[label=\arabic*),leftmargin=2.4em,itemsep=1pt,topsep=2pt]
\item $f_{Z_{\cE}}$ 是格 $\tilde H(S_1,\Z)$ 与 $\tilde H(S_2,\Z)$ 之间的等距；
\item $f$ 的逆映射等于同态
\begin{equation*}
\begin{array}{rccc}
f'\colon & H^*(S_2,\Z) & \longrightarrow & H^*(S_1,\Z)\\[2pt]
& \cup & & \cup\\[2pt]
& \beta & \longmapsto & p_*(Z_{\cE}^{\vee}\cdot\pi^*(\beta)),
\end{array}
\end{equation*}
它由闭链
\begin{equation*}
Z_{\cE}^{\vee}=p^*\sqrt{\td_{S_1}}\cdot\ch(\cE^{\vee})\cdot\pi^*\sqrt{\td_{S_2}}
\end{equation*}
定义，其中 $\cE^{\vee}:=\Rb\cH om(\cE,\cO_{S_1\times S_2})$。
\end{enumerate}
\end{proposition}

% [Tutor Note: 命题 4.1.3 第 1) 条把所构造的等距记作 $f_{Z_{\cE}}$，第 2) 条起改记作 $f$，本页证明与 p.570 也一律用 $f$；原书记号不一致，照录。]

\begin{proof}
与 $\Phi_{\cE}$ 左右伴随的函子彼此同构，它们由公式
\begin{equation*}
\Phi_{\cE}^{*}=\Phi_{\cE}^{!}=\Rf p_*(\cE^{\vee}\otL\pi^*(\,\cdot\,))[2]
\end{equation*}
给出。由于 $\Phi_{\cE}$ 全忠实，复合 $\Phi_{\cE}^{*}\circ\Phi_{\cE}$ 同构于恒等函子 $\id_{\Db(S_1)}$。而恒等函子 $\id_{\Db(S_1)}$ 由对角 $\Delta\subset S_1\times S_1$ 的结构层 $\cO_\Delta$ 定义。
\end{proof}

% [原文 p.570]
由投影公式与 Grothendieck--Riemann--Roch 定理可知，复合 $f'\circ f$ 由闭链
\begin{equation*}
p_1^*\sqrt{\td_{S_1}}\cdot\ch(\cO_\Delta)\cdot p_2^*\sqrt{\td_{S_1}}
\end{equation*}
表示，其中 $p_1$ 与 $p_2$ 是 $S_1\times S_1$ 到两个因子的投影。再用一次 Grothendieck--Riemann--Roch 定理，可知这个闭链等于 $\Delta$。因此复合 $f'\circ f$ 是恒等映射，从而 $f$ 是从 $H^*(S_1,\Z)$ 到 $H^*(S_2,\Z)$ 的同构，因为这两个群都是同秩的自由阿贝尔群。

记 $\nu_S:S\to\Spec\C$ 为 $S$ 的结构态射。于是 $\tilde H(S,\Z)$ 上的配对 $(\alpha,\alpha')$ 可写成 $\nu_*(\alpha^{\vee}\cdot\alpha')$。由投影公式可得
\begin{equation*}
\begin{aligned}
(\alpha,f(\beta))
&=\nu_{S_2,*}\Bigl(\alpha^{\vee}\cdot\pi_*\Bigl(\pi^*\sqrt{\td_{S_2}}\cdot\ch(\cE)\cdot p^*\sqrt{\td_{S_1}}\cdot p^*(\beta)\Bigr)\Bigr)\\
&=\nu_{S_2,*}\pi_*\Bigl(\pi^*(\alpha^{\vee})\cdot p^*(\beta)\cdot\ch(\cE)\cdot\sqrt{\td_{S_1\times S_2}}\Bigr)\\
&=\nu_{S_1\times S_2,*}\Bigl(\pi^*(\alpha^{\vee})\cdot p^*(\beta)\cdot\ch(\cE)\cdot\sqrt{\td_{S_1\times S_2}}\Bigr)
\end{aligned}
\end{equation*}
对任意 $\alpha\in H^*(S_2,\Z)$ 与 $\beta\in H^*(S_1,\Z)$。同法可见
\begin{equation*}
(\beta,f'(\alpha))=\nu_{S_1\times S_2,*}\Bigl(p^*(\beta^{\vee})\cdot\pi^*(\alpha)\cdot\ch(\cE)^{\vee}\cdot\sqrt{\td_{S_1\times S_2}}\Bigr).
\end{equation*}
于是 $(\alpha,f(\beta))=(f'(\alpha),\beta)$。由于 $f'\circ f$ 是恒等映射，可得
\begin{equation*}
(f(\alpha),f(\alpha'))=(f'f(\alpha),\alpha')=(\alpha,\alpha').
\end{equation*}
故 $f$ 是等距。

% [Tutor Note: 上面三行等式的第二个等号把 $\nu_{S_2,*}\pi_*$ 直接换成了 $\nu_{S_1\times S_2,*}$，理由是 $\nu_{S_2}\circ\pi=\nu_{S_1\times S_2}$；原文未写出这一步，属跳步。]

\section{凝聚层导出范畴等价的判别准则}
\takeaway{本节给出全章主定理：两个光滑射影 K3 曲面的凝聚层导出范畴三角等价，当且仅当 Mukai 格之间存在霍奇等距（定理 4.2.1）；再把条件搬到超越闭链格上（定理 4.2.4）。证明分两半：等价必给出霍奇等距，霍奇等距又借 Mukai 的丛模空间造出等价。}

本节给出一个判别准则，判定两个 K3 曲面上凝聚层的导出范畴何时作为三角范畴等价。这个准则的形式很像 K3 曲面的 Torelli 定理：后者断言，两个 K3 曲面 $S_1$ 与 $S_2$ 同构，当且仅当它们的第二上同调格霍奇等距，也就是说，存在等距
\begin{equation*}
H^2(S_1,\Z)\iso H^2(S_2,\Z),
\end{equation*}
它在复上同调上的延拓把 $H^{2,0}(S_1)$ 送到 $H^{2,0}(S_2)$（见 [39]、[27]）。

本章的主要结果如下。

\begin{theorem}{4.2.1}
设 $S_1$ 与 $S_2$ 是复数域 $\C$ 上的两个光滑射影 K3 曲面。则凝聚层的导出范畴 $\Db(S_1)$ 与 $\Db(S_2)$ 作为三角范畴等价，当且仅当 $S_1$ 与 $S_2$ 的 Mukai 格之间存在霍奇等距 $f:\tilde H(S_1,\Z)\iso\tilde H(S_2,\Z)$。
\end{theorem}

这个定理还有另一个版本（定理 4.2.4），也许同样有趣。

我们把定理 4.2.1 的证明拆成两个命题。第一个命题的证明本质上依赖于上一章的主定理 3.2.1，因为它用到一个事实：导出范畴的每个等价都可由积上的对象表示。

% [原文 p.571]
\begin{proposition}{4.2.2}
设 $S_1$ 与 $S_2$ 是两个 K3 曲面，它们的凝聚层导出范畴等价。则超越闭链格 $T_{S_1}$ 与 $T_{S_2}$ 之间存在霍奇等距。
\end{proposition}

\begin{proof}
由定理 3.2.2，积 $S_1\times S_2$ 上存在一个对象 $\cE$，它定义该等价。由命题 4.1.3，$f_{Z_{\cE}}$ 给出 Mukai 格 $\tilde H(S_1,\Z)$ 与 $\tilde H(S_2,\Z)$ 之间的霍奇等距。由于闭链 $Z$ 是代数的，我们得到两个等距
\begin{equation*}
f_{\mathrm{alg}}:-\mathrm{NS}(S_1)\perp\mathrm{U}\iso-\mathrm{NS}(S_2)\perp\mathrm{U}
\qquad\text{与}\qquad
f_\tau:T_{S_1}\iso T_{S_2},
\end{equation*}
其中 $\mathrm{NS}(S_1)$ 与 $\mathrm{NS}(S_2)$ 是 Néron--Severi 格，$T_{S_1}$ 与 $T_{S_2}$ 是超越闭链格。显然 $f_\tau$ 是霍奇等距。
\end{proof}

% [Tutor Note: 此处“由于闭链 $Z$ 是代数的”中的 $Z$ 应指 $Z_{\cE}$——本节先前只定义了 $Z_{\cE}$，未定义过 $Z$；照录原文记号。]

逆命题的证明本质地用到 [31] 的结果，那里研究了 K3 曲面上丛的模空间；同时还用到定理 2.1.5，它给出函子全忠实的判别准则（见 [7]）。

\begin{proposition}{4.2.3}
设 $S_1$ 与 $S_2$ 是两个射影 K3 曲面。假设存在霍奇等距
\begin{equation*}
f:\tilde H(S_2,\Z)\iso\tilde H(S_1,\Z).
\end{equation*}
则凝聚层的有界导出范畴 $\Db(S_1)$ 与 $\Db(S_2)$ 等价。
\end{proposition}

\begin{proof}
令 $v=f(0,0,1)=(r,l,s)$ 与 $u=f(1,0,0)=(p,k,q)$。不妨设 $r>1$。事实上，Mukai 格有两类霍奇等距。第一类由线丛陈特征 $\exp(m)$ 的乘法给出：
\begin{equation*}
\varphi_m(r,l,s)=\Bigl(r,\ l+rm,\ s+(m,l)+\frac{r}{2}m^2\Bigr).
\end{equation*}
第二类是 $r$ 与 $s$ 的互换。利用这两类置换，可以换掉 $f$，使 $r>1$。

% [Tutor Note: $\varphi_m$ 中的 $(m,l)$ 是 $H^2$ 上通常的相交配对，而 p.568 规定 Mukai 格在 $H^2$ 上取负号；原书此式按相交型写出，与前文的符号约定并存，读者套用公式时须留意。]

向量 $v\in\mathrm{U}\perp-\mathrm{NS}(S_1)$ 显然是迷向的，即 $(v,v)=0$。Mukai 在其出色论文 [31] 中证明：此时 K3 曲面 $S_1$ 上存在一个极化 $A$，使得 Mukai 向量与 $v$ 重合、且关于 $A$ 稳定的向量丛的模空间 $\cM_A(v)$ 是光滑射影 K3 曲面。此外，由于存在向量 $u\in\mathrm{U}\perp-\mathrm{NS}(S_1)$ 使 $(v,u)=1$，可知 $\cM_A(v)$ 是精细模空间。于是积 $S_1\times\cM_A(v)$ 上存在泛丛 $\cE$。

\yizhu{这一段的三项条件各司其职，缺一不可：$v$ 迷向保证 $\Ext^1$ 消没，从而不同点的层彼此正交；$v$ 落在正锥内保证模空间非空（这一点原文由 Mukai 的定理一并引用，未展开）；$(v,u)=1$ 给出精细模空间，积上才有泛对象 $\cE$。若只有迷向而没有最后这个本原性条件，得到的只是粗模空间，$\Phi_{\cE}$ 便无从定义。【来源：原书 p.571 及其所引 Mukai [31]（A 级）】}

泛丛 $\cE$ 定义函子 $\Phi_{\cE}:\Db(\cM_A(v))\to\Db(S_1)$。容易看出这个函子满足定理 2.1.5 的条件。确实有 $\Phi_{\cE}(\cO_t)=\cE_t$，其中 $\cE_t$ 是 $S_1$ 上一个稳定丛，且 $v(\cE_t)=v$。所有层 $\cE_t$ 都是单的，当然还有 $\Ext^i(\cE_t,\cE_t)=0$（$i\notin\{0,2\}$）。这给出定理 2.1.5 的条件 2)。

由于 $\cE_t$ 稳定，可得 $\Hom(\cE_{t_1},\cE_{t_2})=0$。由 Serre 对偶，也有 $\Ext^2(\cE_{t_1},\cE_{t_2})=0$。由于向量 $v$ 迷向，还有 $\Ext^1(\cE_{t_1},\cE_{t_2})=0$。

% [原文 p.572]
于是对任意两个不同点 $t_1$ 与 $t_2$，层 $\cE_{t_1}$ 与 $\cE_{t_2}$ 正交。定理 2.1.5 告诉我们：函子 $\Phi_{\cE}$ 全忠实。

事实上 $\Phi_{\cE}$ 不仅全忠实，还是范畴等价。下面的论证基于 [9] 定理 3.3 的证明。记 $\cD$ 为 $\Db(\cM_A(v))$ 在 $\Db(S_1)$ 中的像。由于它是可容许子范畴（见定义 2.2.2），它有右正交与左正交；由于 K3 曲面的典范类平凡，这两个正交重合。于是形如 $\langle\cD^{\perp},\cD\rangle$ 的半正交分解是完全正交的。取 $\cM_A(v)$ 上一个极丰富线丛 $\cL$。所有幂 $\cL^i$ 都是不可分解对象，因而总属于子范畴 $\cD$ 或 $\cD^{\perp}$ 之一；又因为这类对象中任何两个都不完全正交，它们全都属于同一个。而幂 $\{\cL^i\}$ 构成丰富序列（见定义 3.4.1）。由引理 3.4.3，由丰富序列生成的子范畴的正交为 0。于是，既然 $\cD$ 非平凡，就有 $\cD^{\perp}=0$。因此 $\Phi_{\cE}$ 是等价。

其次，由 (47) 定义的闭链 $Z_{\cE}$ 诱导霍奇等距
\begin{equation*}
g:\tilde H(\cM_A(v),\Z)\to\tilde H(S_1,\Z),
\end{equation*}
且 $g(0,0,1)=v=(r,l,s)$。于是 $f^{-1}\circ g$ 也是霍奇等距，并把 $(0,0,1)$ 送到 $(0,0,1)$。因此 $f^{-1}\cdot g$ 诱导 $S_2$ 与 $\cM_A(v)$ 的第二上同调格之间的霍奇等距。故由 Torelli 定理（[39]、[27]），$S_2$ 与 $\cM_A(v)$ 同构。
\end{proof}

% [Tutor Note: 前一句写 $f^{-1}\circ g$，紧接着写 $f^{-1}\cdot g$，同一个复合在两处记号不同；照录原文。]

这个命题与命题 4.1.3 一起证明了定理 4.2.1。定理 4.2.1 还有另一个版本，它用超越闭链格给出导出范畴等价的判别准则。

\begin{theorem}{4.2.4}
设 $S_1$ 与 $S_2$ 是 $\C$ 上两个光滑射影 K3 曲面。则凝聚层的导出范畴 $\Db(S_1)$ 与 $\Db(S_2)$ 作为三角范畴等价，当且仅当 $S_1$ 与 $S_2$ 的超越闭链格之间存在霍奇等距 $f_\tau:T_{S_1}\iso T_{S_2}$。
\end{theorem}

这个断言是定理 4.2.1 与下一个命题的推论。

\begin{proposition}[{{[}33{]}}]{4.2.5}
设 $\phi_1,\phi_2:\mathrm{T}\to\mathrm{H}$ 是格 $\mathrm{T}$ 到偶幺模格 $\mathrm{H}$ 的两个本原嵌入。假设 $\mathrm{H}$ 中的正交补 $\mathrm{N}:=\phi_1(\mathrm{T})^{\perp}$ 以子格的形式包含双曲格 $\mathrm{U}=\begin{pmatrix}0&1\\1&0\end{pmatrix}$。则 $\phi_1$ 与 $\phi_2$ 等价，也就是说，存在 $\mathrm{H}$ 的等距 $\gamma$ 使 $\phi_1=\gamma\phi_2$。
\end{proposition}

事实上，设存在霍奇等距
\begin{equation*}
f_\tau:T_{S_2}\iso T_{S_1}.
\end{equation*}
我们知道，超越闭链格 $T_S$ 在 Mukai 格 $\tilde H(S,\Z)$ 中的正交补同构于格 $-\mathrm{NS}(S)\perp\mathrm{U}$。于是由前一命题（命题 4.2.5），存在等距
\begin{equation*}
f:\tilde H(S_2,\Z)\iso\tilde H(S_1,\Z)
\end{equation*}
使 $f|_{T_{S_2}}=f_\tau$。因此等距 $f$ 也是霍奇等距。故由定理 4.2.1，$S_1$ 与 $S_2$ 上凝聚层的导出范畴等价。

\chapterreview{%
第 4 章在全书中的位置，可以用一句话定位：它是第 3 章表示定理的第一个、也是最完整的应用。第 2 章已经给出了一大批由 flop 相联系的等价簇对，那些例子说明等价确实存在；第 3 章把一切等价归结为积上的一个对象；本章则换一个方向提问——给定两族几何对象，它们的导出范畴何时等价？答案不再来自构造，而来自一个格上的等距条件。与第 2 章相比，这是从举例到分类的转折；与第 5 章相比，本章处理两个不同曲面之间的等价，而第 5 章处理同一簇的自等价群。

知识脉络上有三条线在此汇合。其一是 K3 曲面的经典理论：相交型、Néron--Severi 群、超越闭链格与 Torelli 定理，本来都属于代数几何；其二是 Mukai 关于 K3 曲面上稳定丛模空间的工作，它提供了从格上一个迷向向量造出 K3 曲面的手段；其三是第 3 章的函子—核对象对应，它保证任何等价都由 $S_1\times S_2$ 上的对象给出，从而可以对这个对象取陈特征、造闭链。三条线的接合点就是 §4.1 的闭链 $Z_{\cE}$。

核心结论是一串层层递进的陈述。命题 4.1.3 说：全忠实的 $\Phi_{\cE}$ 诱导 Mukai 格之间的霍奇等距，其证明只是 Grothendieck--Riemann--Roch 与投影公式的反复运用，却是全章的引擎。命题 4.2.2 取正向：等价给出超越闭链格之间的霍奇等距。命题 4.2.3 取反向：先用两类初等等距把 $f$ 归一化成 $r>1$，再把 $(0,0,1)$ 的像实现为某个模空间 $\cM_A(v)$ 上的泛丛，用定理 2.1.5 得全忠实，用丰富序列的消没得本质满，最后借 Torelli 定理把 $\cM_A(v)$ 认作 $S_2$。两半合起来即定理 4.2.1；定理 4.2.4 再把条件搬到超越闭链格上，其格论部分（命题 4.2.5）是纯格论的本原嵌入唯一性。

辨析与提醒有四点。第一，霍奇等距比格等距强：必须把 $H^{2,0}$ 送到 $H^{2,0}$，否则结论不成立，这也正是 Torelli 定理能在证明中出现的原因。第二，Mukai 格里的 $-\mathrm{NS}$ 与 $\mathrm{U}$ 因子并非装饰：负号来自 Serre 对偶，$\mathrm{U}$ 则让 $r$、$s$ 坐标的互换与 $\exp(m)$ 的平移都成为等距，归约 $r>1$ 全靠它。第三，命题 4.1.3 只给出“全忠实 $\Rightarrow$ 等距”，逆向不成立，反向的证明必须动用量空间的几何，这也是本章篇幅的主要去处。第四，全章的基域是 $\C$，Torelli 定理与模空间的射影性都在此使用；换到一般域上，陈述需要重新表述。

评注：本章把“两个 K3 曲面的导出范畴是否等价”化归为一个可以计算的格论问题，这是全书最漂亮的一步。代价是证明中引用了三件外部大定理：Mukai 的模空间定理、K3 曲面的 Torelli 定理、以及命题 4.2.5 的格论结果，作者只给出文献号而不复述证明。Orlov 的贡献在于把这三件工具与第 3 章的表示定理接起来，使范畴等价变成格的算术。也正因如此，本章的结论后来成了检验各种对称性猜想的现成试验场：任何声称在 K3 曲面上给出新等价的理论，都得先过 Mukai 格这一关。}
